\chapter{22-09-26}

hipotesis de rimman 
primo de meisent = P  = 2^r -1 
gimps 

* lema de gauss: un numero p e N, es primo si y solo si tiene h la siguiente propiedad
si p divide un prodcuto ab, entonces p divide a a uno de los dos factores a o b 

6|2.3 y 6 no divide a 2 y 6 no divide a 3 entonces 6 no es primo

si 2| ab entoncs 2|a o 2|b 

\begin{demostracion}[title= demostracion del teorema de euclides?lema de gauss?]
    -> tengo que mostrar que si p es primo entonces p tiene la propiedad * 
    supongamos que p es primo , que p|ab , u wur p no divide a.
    como p no divide a , entonces mod (p,a) = 1 , 
    por el lema de bezout , hay enteros x e y talques xp+ya = 1 
    p | xpb +yab = b 
    si es un numero primo p que no divide a entonces necesariamente tiene que dividir a b 
    <- supongamos que p tiene la propiedade * 
    y mostremos que entonces es primo. supongamos que no es primo. 
    de manera que hay un numero entoer a talquer 1<a<p talque a | p asi que
    hay otro entero b tal que p = a.b 
    ivan ivanivanivanivanandresagustin
    como p tiene la prop * , p divide a a o a b 
    como 1 <a<p p no divide a 
    como p>b>1 , entonces p no divide a b 
    esta contradiccion probino de haber supuesto que p no es primo po lo que tiene que serlo
\end{demostracion}

corolario sea un n un numero primo 
si p divide a un producto a1 a2....an , entonces p divide a algun factor ar 

p1 vale p2 vale pk -> pk+1 vn>2 

\begin{definicion}
    el teorema fundamental de la aritmetica 
    prop : todo elemento n admite una factorizacion como producto de n primos 
    demostracion sea pn = n  es igual a un producto de numeros primos 
    mostremos que pn vale para tod neN por induccion
    p1 vale porque 1 es igual al producto de cero factoreo primos 
    sea n e N , tal que n>=2 y supongamos que as afirmaciones p1,p2, pn-1 valen
    el numero n o es primo o bien es divisible por un numero primo p<n 
    si n es primo entonces n = n es un producto de numeros primos asi que pn vale en ese caso
    supongamos que n no es primo.entonces un numero primo p menor que n u un entero m talque n = pm
    m < n/p < n/2 < n porque p> 2. entonces la hip inductiva que pm vale 
    esto nos dice que hay numeros primos p1,p2,pk tal que m= p1....pk 
    entonces n = pm = p.p1...pk 
    entocnes es un produjcto de primo
    entonces pn vale tambien en este caso
    en cualquier caso pn vale esto completo la induccion
\end{definicion}

\begin{definicion}
    propiedades: si p1,..., pn y q1,... qn son numeros primos y el producto de todos los p1. ....pn = q1 .... qn
    entonces n = m y las listas p1............pn y q1........qm difieren solamente en el orden
    obs: si p y q spm dos primos y p|q , entonces p = q 
\end{definicion}
\begin{demostracion}
    supongamos que p1......pn = q1........qm
    entonces p1 | p1......pn y este producto es igual q1........qm
    asi que p1 divide a alguno de los q-es ,asi que p1 es igual a el 
    sin perdida de generalidad podemos suponer que p1=q1
    entonces p1 p2 .... pn = q1 q2 .. qm = p1 q2 qm 
    asi que p2......pn = q2 ...qm
    asi que p3 ...pn = q3 .... qm 
    entonces  esto lo puedo hacer n veces .¿
    al terminar q1 = p1 , q2 = p2 qm=pn 
    si tacho n terminos me quedan m-n terminos 
    me quedo de un lado 1 = qn+1 qm {m-n terminos}
    qn+1.... qm> 2 * 2 .... 2=2^m-n
    asi que m-n <= 0 asi que m <= n 
\end{demostracion}

\begin{ejemplo}
    12 = 2*2*3 = 2^2 * 3
\end{ejemplo}

i n es un entero , entonces hay una unica forma de elegir un numero natural , numeros primos p1<p2<.....<pr

talques que p1^a1 p2^a2 pn^an (lo mismo que en el ejemplo)

si n= p1^a1 p2^a2 pr^ar p1< --- <pr 

supongamos m es un divisor de n.  m = q1^b1 , ... qs^bs 
